% ---
% Subsección 2.1.7 de los Antecedentes (nueva en TT2; antes 2.1.8 en el borrador del 2026-10-14).
% Justifica codificar un vector de características radiómicas en lugar de los píxeles (D-006, D-013).
% Cada atribución lleva su respaldo verificable con Code/tools/verificar_citas_fichas.py.
% ---

\subsubsection{Representaciones cuánticas de imágenes}
\label{sec:imagenes_cuanticas}

Procesar imágenes en un dispositivo cuántico exige codificar la matriz de píxeles en un estado cuántico. La literatura ha propuesto varias representaciones cuánticas de imágenes (\textit{Quantum Image Representations}, QImRs), cuya viabilidad depende del número de qubits que requieren y del costo del circuito que prepara el estado.

\textbf{Representación FRQI}

% respaldo: «A Flexible Representation of Quantum Images (FRQI) is proposed to provide a representation for images on quantum computers in the form of a normalized state which captures information about colors and their corresponding positions in the images.» (Le2011FRQI, p. 1)
Le et al. introdujeron FRQI (\textit{Flexible Representation of Quantum Images}), que representa una imagen como un estado normalizado que reúne los colores y las posiciones que les corresponden \cite{Le2011FRQI}.
Para una imagen de $2^n \times 2^n$ píxeles, FRQI usa $2n$ qubits para indexar las posiciones $|i\rangle$ y un qubit de color, $2n+1$ qubits en total. La intensidad de cada píxel se codifica en las amplitudes del qubit de color mediante un ángulo $\theta_i \in [0, \pi/2]$ \cite{Le2011FRQI}:
\begin{equation}
  |\psi_{\text{FRQI}}\rangle = \frac{1}{2^n} \sum_{i=0}^{2^{2n}-1} \left( \cos\theta_i |0\rangle + \sin\theta_i |1\rangle \right) \otimes |i\rangle.
\end{equation}

% respaldo: «Theorem 1 (Polynomial preparation theorem) Given a vector θ = (θ0, θ1, . . . , θ22n−1), (n ∈N) of angles, there is a unitary transform P that turns quantum computers from the initialized state, |0⟩⊗2n+1 to the FRQI state» (Le2011FRQI, p. 7)
Los autores enuncian un teorema de preparación: existe una transformación unitaria, construida con compuertas Hadamard y rotaciones controladas, que lleva el estado inicial $|0\rangle^{\otimes 2n+1}$ al estado FRQI \cite{Le2011FRQI}.
% respaldo: «The total number of simple operations used to prepare FRQI state is 2n + 22n × 22n−1 −1 + 22n−1 −2 = 24n −3.22n + 2n.» (Le2011FRQI, p. 7)
Esa preparación usa $2^{4n} - 3 \cdot 2^{2n} + 2n$ operaciones elementales \cite{Le2011FRQI}.
% respaldo: «As proven by the Polynomial Preparation theorem, the total number of simple operations used in the process is polynomial for the number of qubits which are used to encode all positions in an image.» (Le2011FRQI, p. 2)
Los autores la llaman preparación polinomial \cite{Le2011FRQI}. Con $N = 2^{2n}$ píxeles, el término dominante $2^{4n}$ es igual a $N^2$: el costo es cuadrático en el número de píxeles y, por lo tanto, exponencial en el número de qubits de posición.
% respaldo: «For example a 216 position image needs 232 simple gates for preparation.» (Le2011FRQI, p. 8)
Como ejemplo, una imagen de $2^{16}$ posiciones ($256 \times 256$ píxeles) requiere $2^{32}$ compuertas elementales \cite{Le2011FRQI}.

\textbf{Representación NEQR}

% respaldo: «In this paper, based on analysis of existing quantum image representations, a novel enhanced quantum representation (NEQR) for digital images is proposed, which improves the latest flexible representation of quantum images (FRQI).» (Zhang2013NEQR, p. 1)
Zhang et al. propusieron NEQR (\textit{Novel Enhanced Quantum Representation}) como una mejora de FRQI \cite{Zhang2013NEQR}.
% respaldo: «The newly proposed quantum image representation uses the basis state of a qubit sequence to store the gray-scale value of each pixel in the image for the first time, instead of the probability amplitude of a qubit, as in FRQI.» (Zhang2013NEQR, p. 1)
A diferencia de FRQI, NEQR guarda el nivel de gris de cada píxel en un estado de la base computacional de una secuencia de qubits, y no en la amplitud de un solo qubit \cite{Zhang2013NEQR}.
% respaldo: «NEQR needs q + 2n qubits to represent a 2n × 2n image with gray range 2q.» (Zhang2013NEQR, p. 6)
Para una imagen de $2^n \times 2^n$ píxeles con $2^q$ niveles de gris necesita $q + 2n$ qubits: $2n$ para la posición y $q$ para la intensidad \cite{Zhang2013NEQR}. El estado es
\begin{equation}
  |\psi_{\text{NEQR}}\rangle = \frac{1}{2^n} \sum_{Y=0}^{2^n-1} \sum_{X=0}^{2^n-1} |f(Y, X)\rangle |Y\rangle |X\rangle,
\end{equation}
donde $|f(Y, X)\rangle = |C_0 C_1 \cdots C_{q-1}\rangle$ es la codificación binaria del nivel de gris.

% respaldo: «Theorem 2 The whole preparation of NEQR costs no more than O(qn22n) for a 2n × 2n image with gray range 2q. This is equivalent to an approximately quadratic decrease compared to FRQI.» (Zhang2013NEQR, p. 10)
% respaldo: «Theorem 1 The quantum operation UY X, which sets the gray-scale value for its corresponding pixel, costs no more than O(qn) with the assistance of enough ancillary qubits.» (Zhang2013NEQR, p. 10)
Según los autores, la preparación completa de NEQR cuesta a lo sumo $O(q n 2^{2n})$ operaciones, con ayuda de suficientes qubits auxiliares, lo que equivale a una reducción aproximadamente cuadrática respecto a FRQI \cite{Zhang2013NEQR}.
% respaldo: «Through quantum measurements, the original digital image can be accurately retrieved from the NEQR quantum image, not probabilistically as in FRQI.» (Zhang2013NEQR, p. 27)
Además, la imagen original puede recuperarse con exactitud mediante mediciones, y no de forma probabilística como en FRQI \cite{Zhang2013NEQR}.
% respaldo: «It means when the image size is large, the whole preparation will still take unacceptable time consumption.» (Zhang2013NEQR, p. 11)
Aun así, los autores reconocen que, con imágenes grandes, la preparación sigue tomando un tiempo inaceptable \cite{Zhang2013NEQR}. El costo $O(q n 2^{2n})$ es lineal en el número de píxeles, pero sigue siendo exponencial en el número $2n$ de qubits de posición.

\textbf{Comparación de representaciones}

% respaldo: «In this work, we compare and examine the compression properties of four different Quantum Image Representations (QImRs): namely, Tensor Network Representation (TNR), Flexible Representation of Quantum Image (FRQI), Novel Enhanced Quantum Representation NEQR, and Quantum Probability Image Encoding (QPIE).» (Parigi2026QImR, p. 1)
Parigi et al. compararon las propiedades de compresión de cuatro representaciones: TNR (\textit{Tensor Network Representation}), FRQI, NEQR y QPIE (\textit{Quantum Probability Image Encoding}) \cite{Parigi2026QImR}.
% respaldo: «Our simulations show that FRQI and QPIE perform a higher compression of image information than TNR and NEQR.» (Parigi2026QImR, p. 1)
En sus simulaciones, FRQI y QPIE comprimen más la información de la imagen que TNR y NEQR \cite{Parigi2026QImR}.
% respaldo: «In fact, the most off-diagonal elements of FRQI and QPIE Gram matrices are in the range [0.8,1], while for NEQR case they are in the interval [0,0.4]. The TNR Gram matrix shows an intermediate level of compression with respect to FRQI, QPIE, and NEQR, and its elements are in the range [0.4,0.7].» (Parigi2026QImR, p. 7)
La diferencia se ve en las matrices de Gram de los kernels que inducen: la mayoría de los elementos fuera de la diagonal están en $[0.8, 1]$ con FRQI y QPIE, es decir, las imágenes resultan muy parecidas entre sí; en $[0, 0.4]$ con NEQR, y en $[0.4, 0.7]$ con TNR \cite{Parigi2026QImR}.

% respaldo: «Our results indicate that quantum kernels provide comparable classification average accuracy but require exponentially fewer resources for image storage.» (Parigi2026QImR, p. 1)
% respaldo: «Quantum kernels on average perform binary classification with the same accuracy as the classical linear kernel, but with exponentially lower storage cost.» (Parigi2026QImR, p. 7)
Usados en una SVM, estos kernels alcanzaron en promedio la misma exactitud de clasificación binaria que un kernel lineal clásico, con un costo de almacenamiento exponencialmente menor \cite{Parigi2026QImR}.
% respaldo: «However, we highlight that to date the best known protocols for loading an exact representation of the data into an n-qubit state require a number of quantum gates equal to O(2n), which predominates the complexity of quantum algorithms and compromises its potential quantum advantage7.» (Parigi2026QImR, p. 8)
Los autores advierten que los mejores protocolos conocidos para cargar una representación exacta de los datos en un estado de $n$ qubits requieren $O(2^n)$ compuertas, costo que domina la complejidad del algoritmo y compromete su posible ventaja cuántica \cite{Parigi2026QImR}.
% respaldo: «An efficient loading of the classical date into a quantum state is still an open question.» (Parigi2026QImR, p. 8)
La carga eficiente de datos clásicos en un estado cuántico sigue siendo una pregunta abierta \cite{Parigi2026QImR}.

\textbf{Por qué se codifican características radiómicas y no píxeles}

Estas limitaciones explican una decisión de diseño de este trabajo: no codificar la matriz de píxeles de la mamografía, sino un vector reducido de características radiómicas.

% fuente: capítulo 4 (EDA), dimensiones medianas de las mamografías completas y profundidad de 16 bits; el número de qubits y de operaciones es cálculo propio con las expresiones de Zhang et al.
% respaldo: «First, the size of each image is reduced to 16×16 by using the module cv2.INTER_AREA of OpenCV42.» (Parigi2026QImR, p. 5)
Las mamografías completas del CBIS-DDSM miden en mediana unos $5\,037 \times 2\,960$ píxeles, con 16 bits por píxel (capítulo~\ref{sec:eda}). Como estimación propia: para codificarlas con NEQR habría que completarlas a $2^{13} \times 2^{13}$ posiciones ($n = 13$, $q = 16$), lo que exige $16 + 2 \cdot 13 = 42$ qubits y, con la cota de Zhang et al., del orden de $16 \cdot 13 \cdot 2^{26} \approx 1.4 \times 10^{10}$ operaciones de preparación. Eso queda fuera del alcance del hardware actual y de la simulación de vector de estado. Incluso con imágenes reducidas a $16 \times 16$, los experimentos de Parigi et al. muestran kernels de imagen con una exactitud comparable, no superior, a la de un kernel lineal clásico, y el costo de carga sigue siendo el factor dominante.

Por estos motivos, el presente trabajo separa el procesamiento espacial del cómputo cuántico. El recorte de la región de interés, a partir de las máscaras que acompañan al CBIS-DDSM, y la extracción de descriptores de forma, textura e intensidad con PyRadiomics \cite{vanGriethuysen2017} se hacen de forma clásica, y solo el vector de características resultante se codifica en los circuitos cuánticos.
